\properlane{incarnational-marian}{Truth Sprung from Our Earth and the Cloud Overshadowed by the Spirit}
\runninghead{Truth Sprung from Our Earth}

Three details in the Mass of the First Sunday of Advent point directly toward Bethlehem. The opening heading of the 1962 Missal stations the Mass at Rome's great Marian basilica of St.~Mary Major (\latin{Statio ad S.~Mariam maiorem}), the church of the Crib; the Postcommunion prays explicitly for ``the coming solemnities of our reparation'' (\latin{reparationis nostrae ventura solemnia}); and the Communion antiphon sings, ``The Lord will give goodness, and our earth shall yield her fruit'' (Ps.~84:13). When those words are set beside St.~Luke's Gospel, two exegetical questions surface immediately. Why does Christ come \latin{in nube}, ``in a cloud'' (Luke~21:27), and why does Luke alone add ``and all the trees'' (\latin{et omnes arbores}, Luke~21:29) where St.~Matthew mentions only the fig tree (Matt.~24:32)?

Four Fathers and Doctors, joined by liturgical commentators, read these texts in the light of the Incarnation and the Virgin Mary. St.~Ambrose of Milan interprets the cloud of Luke~21:27 as the prophets and the Virgin Mary overshadowed by the Holy Spirit, and the sprouting trees of verse~29 as Israel and all the nations blossoming toward resurrection; St.~Robert Bellarmine expounds Psalm~79:2--4 as begging God to rouse His hidden might and Psalm~84:8 and~13 directly of the Incarnation through the Holy Spirit overshadowing Mary; St.~Augustine identifies Truth sprung from earth in Psalm~84:12 with the Son born of Virgin flesh (\work{Enarrationes in Psalmos}~84.13 and \work{Sermo}~185.1); St.~Bernard of Clairvaux preaches Christ as the Fruit of Psalm~84:13 (\work{Missus est}, Hom.~1) and compares Mary's womb to an unploughed field (\work{Sermones in Adventu}~2); and Sicard of Cremona and Dom Prosper Gu\'eranger read the Communion antiphon of Mary yielding Christ. Read together across the 1962 Mass, their expositions link Christ's coming in majesty directly to His first coming in human flesh: the Word who came in the ``light cloud'' of Mary returns in glory to reap the harvest of Israel and all the nations.

\subsection{The Light Cloud and All the Trees in St.~Ambrose}

St.~Ambrose of Milan (\work{Expositio Evangelii secundum Lucam}~X.37--45) reads Luke~21:25--30 symbolically, taking a very different path from St.~Gregory the Great's literal account of cosmic ruin. Ambrose begins with the celestial signs of verse~25 (X.37--38). The sun that is darkened is Christ Himself, the heavenly Sun, who appears paler or brighter to each viewer according to the measure of faith, and is obscured when widespread unbelief rises like a mist. The moon that ceases to give light is the Church: just as the moon darkens when earth comes between it and the sun, the Church cannot reflect Christ's radiance when carnal vices interpose themselves, or when love of this mortal life shuts out God's light in persecution. The falling stars are the patriarchs, prophets, and saints who shine like luminaries holding fast the word of life; when persecution sharpens, men lose sight of them, and souls wither (\latin{arescentibus}) because the dew of the baptismal font dries up in guilty consciences: \latin{aret perfidia, fides rorat}---``unbelief parches, faith bedews.''

Ambrose next faces a theological puzzle in verse~26: if the heavenly powers (\latin{virtutes caelorum}) already dwell in eternal beatitude, how can they be ``moved'' (X.39--40)? Just as Christ's advent comes to pass both across the whole world and within each individual soul (\latin{in singulis}) that receives Him with all its affections, so the heavenly powers---as the Lord of hosts draws nearer with the fullness of His divinity---are moved by an increase of grace (\latin{augmentum gratiae}), illuminated more brightly as they draw close to Him (Ps.~33:6; 2~Cor.~3:16).

On verse~27, \latin{venientem in nube} (``coming in a cloud,'' X.41--42), Ambrose asks how the Lord who ``hath set his tabernacle in the sun'' (Ps.~18:6) could make a gloomy entrance in storm rain. Scripture instead speaks of clouds moistened with the dew of spiritual grace that temper divine radiance for human eyes. God spoke to Israel in a pillar of cloud through Moses and Joshua; Isaiah and Ezekiel are clouds flying above us, revealing the holiness of the Trinity through cherubim and seraphim (Isa.~60:8)---the angelic praise echoed in this Sunday's Preface of the Most Holy Trinity, which confesses the co-equal divinity of the Father, the only-begotten Son, and the Holy Spirit. Every prophet is a cloud in which Christ comes, and above all the Virgin Mary: ``He came also in a light cloud, incarnate of the Virgin; for the prophet saw Him coming like a cloud from the East'' (citing Isa.~19:1). Why did Isaiah call that cloud \emph{light}? Because ``no earthly vices weighed her down.'' Ambrose confirms the identification from Luke's own Infancy narrative: ``Look upon the cloud on which the Holy Spirit came down, and the power of the Most High overshadowed'' (Luke~1:35).

Finally, Ambrose turns to the textual difference between Matthew~24:32 (``the fig tree'') and Luke~21:29 (``the fig tree, and all the trees,'' X.43--45). Where Gregory saw the world's ruin as its ironic fruit, Ambrose interprets the sprouting branches as genuine fruitfulness: when the fruit grows green on all the trees and the fig tree itself blossoms, ``so that every tongue may confess to God, the Jewish people also confessing,'' then we look for the Lord's coming, when the fruits of the resurrection will be harvested as in summer. For Ambrose, the fig tree signifies Israel and ``all the trees'' signify all the nations of the Church blossoming with faith. He does not erase the Gospel's warning, however: the fig tree carries a double figure (\latin{duplicem figuram}). Beside fruitful faith stands the barren fig tree covered only in leaves---the fragile boasting of sinners and Antichrist sewing fig leaves together to hide guilty consciences, ``the garments of exiles from paradise'' (Gen.~3:7).

\subsection{Our Earth Yielding Her Fruit}

Where Ambrose reads the Gospel's cloud of the Virgin overshadowed by the Spirit, St.~Robert Bellarmine, St.~Augustine, and St.~Bernard of Clairvaux trace the same mystery of the Incarnation in the Mass's psalm chants and prayers. At St.~Mary Major, the Introit, Gradual, and Offertory from Psalm~24 voice the waiting of Israel and the whole Church (\latin{totus Christus}, ``the whole Christ,'' Augustine \work{Enarr.~in~Ps.}~24.1) for the promised Savior. The Collect (\latin{Excita, quaesumus}), whose \latin{Qui vivis} ending addresses the eternal Son directly, echoes Psalm~79:3 (\latin{Excita potentiam tuam, et veni}), which Bellarmine (\work{Explanatio} on Ps.~79:2--4) expounds as begging God to rouse His seemingly hidden might and come to save His people. At the Alleluia (Ps.~84:8, ``Show us, O Lord, thy mercy; and grant us thy salvation''), Bellarmine (\work{Explanatio} on Ps.~84:8) expounds \latin{salutare tuum} directly of the Incarnate Son: ``Send us your Son for a Savior,'' fulfilling St.~Paul's words that ``the grace of God hath appeared to all men'' and ``the goodness and kindness of our Savior God appeared'' (Titus~2:11, 3:4).

At the Communion (Ps.~84:13, ``The Lord will give goodness: and our earth shall yield her fruit''), Bellarmine asks how barren human soil can bear saving fruit and answers from the Annunciation (\work{Explanatio} on Ps.~84:13): truth springs from the earth like wild flowers nourished solely by sunlight and rain from heaven rather than crops sown by human hands. ``For the Lord will give goodness'' means that God ``will send his Holy Spirit from heaven, who will overshadow a virgin, and thus our land, which was never ploughed nor sown, and was altogether an untouched virgin, will yield her fruit''---Christ, who calls Himself ``the flower of the field, and the lily of the valleys'' (Cant.~2:1).

St.~Augustine grounds that Incarnational reading in the immediately preceding verse, Psalm~84:12 (\latin{Veritas de terra orta est}, ``Truth hath sprung out of the earth''). In \work{Enarrationes in Psalmos}~84.13 (Coxe~85.10), he asks: ``What is truth? The Son of God. What is the earth? Flesh.'' Though the Word made heaven and earth before time, He was born of the Virgin Mary so that righteousness might look down from heaven and justify sinners. In his Christmas Day sermon \work{Sermo}~185.1 (\latin{In Natali Domini}~II, PL~38, cols.~997--998), Augustine unfolds the same verse before issuing a wake-up call that echoes Advent's summons: ``Truth which is in the bosom of the Father sprang from the earth so as to be also in a mother's bosom \dots\ Truth for whom heaven does not suffice sprang from the earth so as to be laid in a manger.'' Asking for whose benefit such majesty stooped so low---none of His own, but ours---he cries out, ``Awake, O man: for you God became man. Rise, you who sleep, and arise from the dead, and Christ shall enlighten you'' (Eph.~5:14), and closes with the warning, ``You would have perished, had He not come.'' William Durandus (\work{Rationale}~VI.3, fol.~261r) draws the same exchange from the Epistle's command to put on Christ (Rom.~13:14a)---``for this reason the Son of God became man, that man might become a son of God''---and notes (\work{Rationale}~VI.2, fol.~256r) that the angels' Bethlehem hymn \latin{Gloria in excelsis} is kept silent in Advent so it may be sung more devoutly on Christmas night.

St.~Bernard of Clairvaux (\work{Missus est}, Hom.~1, Browne \& Nolan pp.~56, 58--60, esp.~p.~59) traces Psalm~84:13 across salvation history. The promises given to Abraham, Isaac, and Jacob are the seed cast upon the earth; the wonders of the Exodus, the visions of the prophets, and the institutions of kingdom and priesthood are the blossom of Nazareth; and Christ Himself is the Fruit, for which Bernard quotes Psalm~84:13 (``The Lord will give goodness, and our earth shall yield her Fruit'') beside God's promise to David in Psalm~131:11 (``Of the Fruit of thy womb I will set upon thy throne''). Just as the flower falls once the fruit appears, the Old Testament figures passed away when Truth appeared in human flesh. In his \work{Second Sermon for Advent} (p.~19), Bernard draws the Marian image directly on Isaiah~11:1 and Canticle~2:1, comparing Mary's inviolate womb to an unploughed, unsown field yielding the divine Flower. Liturgical commentators apply Psalm~84:13 along the same Marian line: Sicard of Cremona (\work{Mitrale}~V.2, PL~213, col.~205) writes that by the goodness of the Holy Spirit ``our earth,'' the Blessed Virgin, conceives and yields her fruit, God and man; Dom Prosper Gu\'eranger (\work{The Liturgical Year}~I, p.~137) likewise identifies ``our Earth'' with the Blessed Virgin ``made fruitful by the dew of heaven'' after Isaiah, adding that the divine Fruit has just given Himself by anticipation to the faithful in Communion. Having asked in the Secret (\latin{Haec sacra nos}) to come purer to the Mysteries' divine Source (\latin{principium}), the congregation gathered at St.~Mary Major prays in the Postcommunion (\latin{Suscipiamus}) for \latin{reparationis nostrae ventura solemnia}, asking that mercy received in the temple prepare us to greet the approaching festival of the Nativity with fitting honor.

\subsection{Tensions and Limits of This Perspective}

Ambrose himself guards his interpretation against easy optimism: his double figure of Luke~21:29 warns that the same fig tree can stand for barren leaves of deceit rather than genuine fruit, and his opening paragraphs describe unbelief parching baptismal grace into spiritual drought. Its chief exegetical limit lies on the surface of St.~Luke's text: phrases such as ``the roaring of the sea and of the waves'' and ``men withering away for fear'' carry a literal gravity of public judgment that Gregory and Cyril confront directly, whereas Ambrose's treatment works predominantly through christological and ecclesial allegory.

\subsection{The Four Senses}

\begin{fourSenses}
\item[Literal.] Israel waits on God and pleads for His saving mercy and fruitful blessing (Pss.~24, 84); St.~Paul calls believers to live in the newly dawned day of Christ (Rom.~13:11--14a); Jesus foretells the coming of the Son of Man in a cloud and the blossoming of the fig tree and all trees (Luke~21:25--33); the Roman Church gathers at St.~Mary Major to prepare for the coming solemnities of our reparation (see \hyperref[sec:date-location]{\work{Scriptural Date and Location}} for scriptural chronology).
\item[Allegorical.] The \latin{salutare} of Psalm~84:8 is Christ incarnate (Augustine, Bellarmine); the \latin{nube} of Luke~21:27 signifies the prophets and the Virgin Mary overshadowed by the Holy Spirit (Ambrose, \work{Exp.~in~Luc.}~X.42); \latin{benignitas} is the Holy Spirit and \latin{terra nostra} is the untouched Virgin Mary yielding Christ (Bellarmine, Sicard), while Bernard names Christ the Fruit of Psalm~84:13 (\work{Missus est}, Hom.~1, p.~59) and compares Mary's womb to an unploughed field (\work{In~Adventu}~2, p.~19); the fig tree and ``all the trees'' (\latin{et omnes arbores}) signify Israel and all Gentile nations blossoming into resurrection faith (Ambrose, X.44).
\item[Moral.] Receiving Christ's coming within each individual soul (\latin{in singulis}, Ambrose~X.39), guarding the dew of baptismal faith against the parching drought of unbelief (\latin{fides rorat}, X.38), putting on Christ in daily life (Rom.~13:14a), and preparing with reverent purity for Christmas.
\item[Anagogical.] The summer harvest of the general resurrection when the seed planted at the Incarnation reaches its full flowering across every nation in God's kingdom, united with the angels and seraphim of the Trinity Preface in eternal adoration of the Triune God.
\end{fourSenses}
